% ================================================================
\section{The Riemann Sphere: Geometric Foundations}
\label{sec:riemann-sphere}
% ================================================================

The Zeta subsystem (\lean{Zeta/*.lean}, 10~files, $\sim$3{,}400~lines,
zero sorry except 1 structural axiom) establishes a complete
geometric interpretation of the critical line, the functional equation,
and the zero distribution on the Riemann sphere~$S^2$.

\subsection{The Three Realities}

The Mirror Geometry (\lean{MirrorGeometry.lean}, 23~theorems, 0~sorry)
partitions~$\CC$ into three sectors separated by the critical line:

\begin{center}
\small
\begin{tabular}{@{}lll@{}}
\toprule
\textbf{Sector} & \textbf{Definition} & \textbf{Physical Character} \\
\midrule
Positive Reality & $\re s > \tfrac{1}{2}$ & Euler product converges; $\zeta \neq 0$ for $\re s > 1$ \\
Zero Dimension   & $\re s = \tfrac{1}{2}$ & Critical line: all nontrivial zeros live here (RH) \\
Negative Reality & $\re s < \tfrac{1}{2}$ & Trivial zeros at $s = -2, -4, \ldots$ \\
\bottomrule
\end{tabular}
\end{center}

\begin{principle}[The Mirror Map]
The functional equation acts as an involution
$s \mapsto 1 - s$ (\lean{mirror\_involution}), swapping Positive
and Negative Reality (\lean{mirror\_swaps\_sectors}). The critical line
$\re s = \tfrac{1}{2}$ is the \emph{fixed locus} of this mirror
(\lean{mirror\_fixes\_critical\_line}).

On the critical line, the mirror equals conjugation
(\lean{mirror\_eq\_conj\_on\_critical\_line}):
\[
  1 - \bigl(\tfrac{1}{2} + it\bigr) = \tfrac{1}{2} - it
  = \overline{\tfrac{1}{2} + it}
\]
This coincidence is the geometric reason why $\Lambda_0(\tfrac{1}{2}+it) \in \RR$
(\lean{critical\_line\_reality}), reducing the zero-finding problem
from two real dimensions to one.
\end{principle}

\subsection{Centered Coordinates and the Great Circle}

The Great Circle (\lean{RiemannSphere.lean}, 23~theorems, 0~sorry)
introduces the centered coordinate $w = s - \tfrac{1}{2}$ (\lean{centered}),
under which:
\begin{itemize}
  \item The critical line $\re s = \tfrac{1}{2}$ becomes the
    \emph{imaginary axis} $\re w = 0$
    (\lean{critical\_line\_is\_imaginary\_axis}).
  \item The functional equation $s \mapsto 1-s$ becomes
    \emph{negation} $w \mapsto -w$
    (\lean{func\_eq\_is\_negation}).
  \item The completed zeta function $\Lambda_0$ becomes an
    \emph{even function} of~$w$:
    $\Lambda_0(\tfrac{1}{2}+w) = \Lambda_0(\tfrac{1}{2}-w)$
    (\lean{completedZeta\_even\_in\_w}).
\end{itemize}

Under stereographic projection, lines through the origin of~$\CC$
map to great circles on~$S^2$. Since the critical line \emph{is}
the imaginary axis in centered coordinates---a line through the
origin---it maps to a \emph{great circle}:

\begin{theorem}[The Great Circle Theorem, \lean{great\_circle\_zero\_structure}]
\label{thm:great-circle}
The critical line is a great circle on the Riemann sphere.
The zeros of~$\Lambda_0$ on this great circle form antipodal pairs:
if $w_0 = t_0 \cdot i$ is a zero, then $-w_0 = -t_0 \cdot i$ is also
a zero (\lean{zeros\_are\_antipodal}).
The great circle divides~$S^2$ into two hemispheres:
\begin{itemize}
  \item \textbf{Eastern Hemisphere}: $\re w > 0$ (Positive Reality, $\re s > \tfrac{1}{2}$)
  \item \textbf{Western Hemisphere}: $\re w < 0$ (Negative Reality, $\re s < \tfrac{1}{2}$)
\end{itemize}
The functional equation ($w \mapsto -w$) is a rotation by~$\pi$
swapping East and West (\lean{negation\_swaps\_hemispheres}).
\end{theorem}

\begin{correspondence}[Geodesic Principle]
In general relativity, geodesics are the paths of zero net force.
The critical line, being a great circle, is the \emph{geodesic}
of the Riemann sphere's geometry. The RH asserts that
all nontrivial zeros of~$\zeta$ sit on this geodesic---they
follow the path of least resistance through the arithmetic
landscape. This is the geometric incarnation of ``the primes
choose the simplest possible structure.''
\end{correspondence}

\subsection{The Klein Four-Group and Its Degeneration}

The Four-Fold Symmetry (\lean{FourFoldSymmetry.lean}, 36~theorems, 0~sorry)
formalizes the Klein four-group $V_4 = \{e, M, C, MC\}$ acting on~$\CC$:
\[
  e : s \mapsto s, \qquad
  M : s \mapsto 1-s, \qquad
  C : s \mapsto \bar{s}, \qquad
  MC : s \mapsto 1-\bar{s}
\]
with multiplication table and fixed-point analysis. In centered
coordinates, this simplifies to:
\[
  e : w \mapsto w, \qquad
  M : w \mapsto -w, \qquad
  C : w \mapsto \bar{w}, \qquad
  MC : w \mapsto -\bar{w}
\]

On the critical line ($w = ti$), the Klein group \textbf{degenerates}
(\lean{negation\_eq\_conj\_on\_great\_circle}): $-w = -(ti) = \overline{ti} = \bar{w}$,
so the mirror $M$ and conjugation~$C$ \emph{coincide}. The group
$V_4 \cong \ZZ/2 \times \ZZ/2$ collapses to $\ZZ/2$.

The Circle Quadruplet (\lean{CircleQuadruplet.lean}, 12~theorems, 0~sorry)
extends this to the action of $V_4$ on circles, showing that the
Klein group maps the unit circle at height~$t$ to four related circles,
which degenerate at the equator ($t = 0$).

\begin{remark}[Physical interpretation]
The Klein degeneration is the number-theoretic analogue of
\emph{symmetry breaking} in gauge theory. The full $V_4$ symmetry
exists everywhere, but on the critical line---the equator---two
independent symmetries merge into one. This ``enhanced symmetry'' at
the equator is precisely what forces $\Lambda_0$ to be real there,
collapsing the 2D zero-finding problem to 1D.
\end{remark}

\subsection{The Algebraic Tower}

Three files extend the geometric picture into algebraic number theory:

\begin{itemize}
  \item \textbf{The Kummer Tower} (\lean{KummerTower.lean}, 0~sorry):
    extends the Cayley-Dickson construction (reals $\to$ complex $\to$
    quaternions $\to$ octonions $\to \ldots$) into the $p$-adic direction
    via Kummer extensions $\mathbb{Q}(\zeta_p)$. Each prime~$p$ contributes a
    ``floor'' to the tower, and the Euler product is the product of
    all floors.

  \item \textbf{The Spectral Tower} (\lean{SpectralTower.lean}, 0~sorry):
    formalizes the Fourier harmonics in the imaginary direction of the
    critical strip, connecting Fourier analysis on $\re s = \sigma$
    to the spectral decomposition of~$\zeta$.

  \item \textbf{Tower Fusion} (\lean{TowerFusion.lean}):
    states the \emph{Rigidity Axiom}: the assertion that
    all nontrivial zeros of~$\zeta$ in the critical strip have
    $\re s = \tfrac{1}{2}$.  This is precisely the Riemann Hypothesis,
    stated as a single axiom.  The key result:
    \lean{F1\_dream} (in \lean{MirrorGeometry.lean}) is \emph{graduated}
    from an independent axiom to a theorem derived from \lean{tower\_fusion}.
\end{itemize}

\begin{remark}[The Dream of $\mathbb{F}_1$]
The ``field with one element'' $\mathbb{F}_1$, if formalized, would reduce
$\operatorname{Spec}(\ZZ)$ (an infinite-genus surface) to~$S^2$ (a sphere), on which
the Riemann Hypothesis is trivially true by the Weil conjectures.
\lean{F1\_dream} records this structural statement, now derived from
the Tower Fusion axiom rather than stated independently.
\end{remark}

\subsection{Theorem Census}

\begin{center}
\small
\begin{tabular}{@{}llrl@{}}
\toprule
\textbf{File} & \textbf{Content} & \textbf{Theorems} & \textbf{Sorry} \\
\midrule
\lean{MirrorGeometry.lean}   & Three Realities, Mirror, S-Duality   & 14 & 0 \\
\lean{SilenceAndEcho.lean}   & Trivial zero duality                 & 9  & 0 \\
\lean{FourFoldSymmetry.lean} & Klein $V_4$ group, Four Noble Zeros   & 36 & 0 \\
\lean{StripGeometry.lean}    & Circle-strip geometry, Hardy~$Z$      & 14 & 0 \\
\lean{HardyZFunction.lean}   & Hardy $Z$-function, ring contraction  & 13 & 0 \\
\lean{RiemannSphere.lean}    & Great Circle, hemisphere bisection    & 23 & 0 \\
\lean{CircleQuadruplet.lean} & $V_4$ action on circles               & 12 & 0 \\
\lean{KummerTower.lean}      & $p$-adic tower extension              & 11 & 0 \\
\lean{SpectralTower.lean}    & Fourier harmonics                     & 10 & 0 \\
\lean{TowerFusion.lean}      & Rigidity Axiom ($\equiv$ RH)          & 8  & 0 \\
\midrule
\textbf{Total}               &                                       & \textbf{150} & \textbf{0} \\
\bottomrule
\end{tabular}
\end{center}
